% DOCUMENT CLASS AND CODING
\usepackage[utf8]{inputenc}
\usepackage[T1]{fontenc}
% Dejar que la clase controle márgenes y estilos; añadir rutas de imágenes
\graphicspath{{images/}{reference/}{.}}

% PACKAGES (solo paquetes complementarios; la clase oficial gestiona layout/chapters/pagestyles)
\usepackage{graphicx}
\usepackage{float}
\usepackage{xcolor}
\usepackage{lmodern}
\usepackage{microtype}
\usepackage[hidelinks]{hyperref}
% Acronyms (print only used acronyms)
\usepackage[printonlyused]{acronym}
\usepackage{listings}
\usepackage{booktabs}
% Tablas largas
\usepackage{longtable}
\usepackage{amsmath}
\usepackage{subcaption}

%%%%%%% COLORES %%%%%%%
\definecolor{myblue}{RGB}{0,82,155}
% Colores para código
\definecolor{codegreen}{rgb}{0,0.6,0}
\definecolor{codegray}{rgb}{0.5,0.5,0.5}
\definecolor{codepurple}{rgb}{0.58,0,0.82}
\definecolor{backcolour}{rgb}{0.95,0.95,0.92}

%%%%%%% CONFIGURACIÓN DE CÓDIGO %%%%%%%

% Estilo General por defecto
\lstdefinestyle{mystyle}{
    backgroundcolor=\color{backcolour},   
    commentstyle=\color{codegreen},
    keywordstyle=\color{magenta},
    numberstyle=\tiny\color{codegray},
    stringstyle=\color{codepurple},
    basicstyle=\ttfamily\small,
    breakatwhitespace=false,         
    breaklines=true,                 
    captionpos=b,                    
    keepspaces=true,                 
    numbers=left,                    
    numbersep=5pt,                  
    showspaces=false,                
    showstringspaces=false,
    showtabs=false,                  
    tabsize=2,
    frame=single, % Añade un borde al código
    rulecolor=\color{gray!30}
}

% Aplicar estilo general
\lstset{style=mystyle}

% --- LENGUAJES ESPECÍFICOS ---

% SQL
\lstdefinestyle{sql}{
    language=SQL,
    style=mystyle,
    keywordstyle=\color{blue}\bfseries,
    stringstyle=\color{red}
}

% JAVA
\lstdefinestyle{java}{
    language=Java,
    style=mystyle,
    keywordstyle=\color{myblue}\bfseries,
}

% XML
\lstdefinestyle{xml}{
    language=XML,
    style=mystyle,
    morestring=[b]",
    morestring=[s]{>}{<},
    morecomment=[s]{<?}{?>},
    stringstyle=\color{black},
    identifierstyle=\color{blue},
    keywordstyle=\color{cyan},
    morekeywords={xmlns,version,type}
}

% JSON
\lstdefinelanguage{json}{
    basicstyle=\ttfamily\small,
    numbers=left,
    numberstyle=\tiny\color{codegray},
    stepnumber=1,
    numbersep=8pt,
    showstringspaces=false,
    breaklines=true,
    frame=single,
    backgroundcolor=\color{backcolour},
    % Aquí está la magia: Mapeamos colores Y caracteres especiales
    literate=
     *{0}{{{\color{blue}0}}}{1}
      {1}{{{\color{blue}1}}}{1}
      {2}{{{\color{blue}2}}}{1}
      {3}{{{\color{blue}3}}}{1}
      {4}{{{\color{blue}4}}}{1}
      {5}{{{\color{blue}5}}}{1}
      {6}{{{\color{blue}6}}}{1}
      {7}{{{\color{blue}7}}}{1}
      {8}{{{\color{blue}8}}}{1}
      {9}{{{\color{blue}9}}}{1}
      {:}{{{\color{black}{:}}}}{1}
      {,}{{{\color{black}{,}}}}{1}
      {\{}{{{\color{black}{\{}}}}{1}
      {\}}{{{\color{black}{\}}}}}{1}
      {[}{{{\color{black}{[}}}}{1}
      {]}{{{\color{black}{]}}}}{1}
      {á}{{\'a}}1
      {é}{{\'e}}1
      {í}{{\'i}}1
      {ó}{{\'o}}1
      {ú}{{\'u}}1
      {Á}{{\'A}}1
      {É}{{\'E}}1
      {Í}{{\'I}}1
      {Ó}{{\'O}}1
      {Ú}{{\'U}}1
      {ñ}{{\~n}}1
      {Ñ}{{\~N}}1,
}